% Sections won't have numbers
% Make the page area as large as possible
% Set the first level numbering to arabic and the second level
% to alphabetic. The remaining two levels are arabic
%Change to \pagestyle{empty} to suppress numbers on following pages
%\pagestyle{empty}
% Do not reset outermost enumerate counter


\documentclass{article}
%%%%%%%%%%%%%%%%%%%%%%%%%%%%%%%%%%%%%%%%%%%%%%%%%%%%%%%%%%%%%%%%%%%%%%%%%%%%%%%%%%%%%%%%%%%%%%%%%%%%%%%%%%%%%%%%%%%%%%%%%%%%%%%%%%%%%%%%%%%%%%%%%%%%%%%%%%%%%%%%%%%%%%%%%%%%%%%%%%%%%%%%%%%%%%%%%%%%%%%%%%%%%%%%%%%%%%%%%%%%%%%%%%%%%%%%%%%%%%%%%%%%%%%%%%%%
\usepackage{amsmath}
\usepackage{sw20exm2}
\usepackage[compat2]{geometry}

\setcounter{MaxMatrixCols}{10}
%TCIDATA{OutputFilter=LATEX.DLL}
%TCIDATA{Version=5.50.0.2953}
%TCIDATA{<META NAME="SaveForMode" CONTENT="1">}
%TCIDATA{BibliographyScheme=Manual}
%TCIDATA{Created=Monday, May 21, 2012 18:55:59}
%TCIDATA{LastRevised=Friday, November 05, 2021 09:57:52}
%TCIDATA{<META NAME="ViewSettings" CONTENT="0">}
%TCIDATA{<META NAME="GraphicsSave" CONTENT="32">}
%TCIDATA{<META NAME="DocumentShell" CONTENT="Exams and Syllabi\SW\SW Exam #1 - 8.5x14 Page">}
%TCIDATA{CSTFile=Exam.cst}

\setlength{\topmargin}{-0.75in}
\setlength{\textheight}{12.25in}
\setlength{\oddsidemargin}{0.0in}
\setlength{\evensidemargin}{0.0in}
\setlength{\textwidth}{6.5in}
\def\labelenumi{\arabic{enumi}.}
\def\theenumi{\arabic{enumi}}
\def\labelenumii{(\alph{enumii})}
\def\theenumii{\alph{enumii}}
\def\p@enumii{\theenumi.}
\def\labelenumiii{\arabic{enumiii}.}
\def\theenumiii{\arabic{enumiii}}
\def\p@enumiii{(\theenumi)(\theenumii)}
\def\labelenumiv{\arabic{enumiv}.}
\def\theenumiv{\arabic{enumiv}}
\def\p@enumiv{\p@enumiii.\theenumiii}
\pagestyle{plain}
\setcounter{secnumdepth}{0}
\newtheorem{theorem}{Theorem}
\newtheorem{acknowledgement}[theorem]{Acknowledgement}
\newtheorem{algorithm}[theorem]{Algorithm}
\newtheorem{axiom}[theorem]{Axiom}
\newtheorem{case}[theorem]{Case}
\newtheorem{claim}[theorem]{Claim}
\newtheorem{conclusion}[theorem]{Conclusion}
\newtheorem{condition}[theorem]{Condition}
\newtheorem{conjecture}[theorem]{Conjecture}
\newtheorem{corollary}[theorem]{Corollary}
\newtheorem{criterion}[theorem]{Criterion}
\newtheorem{definition}[theorem]{Definition}
\newtheorem{example}[theorem]{Example}
\newtheorem{exercise}[theorem]{Exercise}
\newtheorem{lemma}[theorem]{Lemma}
\newtheorem{notation}[theorem]{Notation}
\newtheorem{problem}[theorem]{Problem}
\newtheorem{proposition}[theorem]{Proposition}
\newtheorem{remark}[theorem]{Remark}
\newtheorem{solution}[theorem]{Solution}
\newtheorem{summary}[theorem]{Summary}
\newenvironment{proof}[1][Proof]{\noindent\textbf{#1.} }{\ \rule{0.5em}{0.5em}}
\input{tcilatex}
\begin{document}


\begin{center}
\begin{equation*}
\FRAME{itbpF}{1.1052in}{0.8614in}{0in}{}{}{Figure}{\special{language
"Scientific Word";type "GRAPHIC";maintain-aspect-ratio TRUE;display
"USEDEF";valid_file "T";width 1.1052in;height 0.8614in;depth
0in;original-width 2.6809in;original-height 2.0833in;cropleft "0";croptop
"1";cropright "1";cropbottom "0";tempfilename
'R23O0G00.bmp';tempfile-properties "XPR";}}
\end{equation*}%
\textbf{Electromagnetismo}

\emph{Gu\'{\i}a N}$%
%TCIMACRO{\U{b0}}%
%BeginExpansion
{{}^\circ}%
%EndExpansion
7$\emph{\ - Capacitancia y Diel\'{e}ctricos}

\textit{Profesor: Arturo Fern\'{a}ndez P\'{e}rez}
\end{center}

\begin{enumerate}
\item Un capacitor de placas paralelas tiene placas circulares de $8,22$ $%
[cm]$ de radio y $1,31$ $[cm]$ de separaci\'{o}n \textquestiondown Cu\'{a}l
es su capacitancia con aire entre sus placas? \textquestiondown Y con tefl%
\'{o}n ($\kappa =2,1$) entre ellas? \textquestiondown Qu\'{e} carga aparece
en los electrodos, en ambos casos, si se aplica una diferencia de potencia
de $116$ $[V]$? \textbf{R:} $C_{0}=14,34$ [pF], $C_{1}=30,11$ [pF], $%
q_{0}=1,66$ [nC], $q_{1}=3,49$ [nC]

\item Un capacitor esf\'{e}rico de $2$ $[\mu F]$ est\'{a} compuesto de dos
esferas met\'{a}licas, una con radio dos veces mayor que la otra. Si la regi%
\'{o}n entre las esferas es aire, determine el volumen de esta regi\'{o}n. 
\textbf{R:} $2,13\times 10^{13}$ $[m^{3}]$

\item Un capacitor cil\'{\i}ndrico tiene conductores exterior e interior
cuyos radios est\'{a}n en proporci\'{o}n $4:1$. El conductor interior se va
a sustituir por un alambre cuyo radio es la mitad del radio del conductor
interior original. \textquestiondown En qu\'{e} factor debe incrementarse la
longitud para obtener una capacitancia igual a la del condensador original? 
\textbf{R: }En 1,5 veces.

\item Un capacitor de placas paralelas de $16$ $[pF]$ se carga por medio de
una bater\'{\i}a de $10$ $[V]$. Si cada placa del capacitor tiene un \'{a}%
rea de $5$ $[cm^{2}]$

\begin{enumerate}
\item \textquestiondown Cu\'{a}l es el valor de la energ\'{\i}a almacenada
en el capacitor? \textbf{R: }$8\times 10^{-10}$ [J]

\item Si se inserta una l\'{a}mina de PVC ($\kappa =3,18$) entre las placas
del capacitor \textquestiondown Cu\'{a}nta energ\'{\i}a puede almacenar? 
\textbf{R: }$2,544\times 10^{-9}$ [J]
\end{enumerate}

\item Un banco de 2100 capacitores de 5 $[\mu F]$ conectados en paralelo se
usa para almacenar energ\'{\i}a el\'{e}ctrica \textquestiondown Cu\'{a}nta
carga se puede almacenar a $55000$ [$V]$? \textbf{R: }$577,5$ [C]

\item Dos capacitores de $2,12$ [$\mu F]$ y $3,88$ [$\mu F],$ est\'{a}n
conectados en serie por una diferencia de potencial de $328$ [$V]$.\ Calcule
la energ\'{\i}a total almacenada. \textbf{R: }$0,0736$ [J]

\item Un capacitor de placas paralelas llenas de aire tiene una capacitancia
de $1,32$ [$pF]$. La separaci\'{o}n de las placas se duplica y entre ellas
se inserta cera.\ La nueva capacitancia es $2,57$ [$pF]$. Determine la
constante diel\'{e}ctrica de la cera. \textbf{R: }$\kappa =3,89$

\item Se construye un capacitor (Fig. 1) con dos placas cuadradas de $12,0$ $%
[cm]$ de lado, separados una distancia de $4,5$ $[mm]$. La mitad del espacio
se rellena con \textit{Plexiglas} ($\kappa =3,4$) y la otra con aire. Una
bater\'{\i}a de $18$ [$V]$ se conecta al condensador \textquestiondown Cu%
\'{a}l es la capacitancia de este dispositivo? \textbf{R: }$62,30$ [pF] 
\textquestiondown Cu\'{a}nta energ\'{\i}a puede almacenar? \textbf{R: }$%
10,09\times 10^{-9}$ [J] Si se remueve el \textit{Plexiglas} 
\textquestiondown cu\'{a}nta energ\'{\i}a se puede almacenar? \textbf{R: }$%
4,58\times 10^{-9}$ [J]%
\begin{equation*}
\FRAME{itbpFU}{1.9882in}{0.8864in}{0in}{\Qcb{Figura 1}}{}{Figure}{\special%
{language "Scientific Word";type "GRAPHIC";maintain-aspect-ratio
TRUE;display "USEDEF";valid_file "T";width 1.9882in;height 0.8864in;depth
0in;original-width 6.5855in;original-height 2.9196in;cropleft "0";croptop
"1";cropright "1";cropbottom "0";tempfilename
'R23O0G01.wmf';tempfile-properties "XPR";}}
\end{equation*}%
\pagebreak

\item Para el circuito de la Fig. 2, donde $C_{1}=12,5$ $[\mu F]$, $%
C_{2}=5,30$ $[\mu F]$, $C_{3}=4,50$ $[\mu F]$

\begin{enumerate}
\item Determine la capacitancia equivalente. \textbf{R:} $C_{eq}=3,571$ [$%
\mu F$]

\item Si $V=12,5$ $[V]$, determine la carga en cada condensador. \textbf{R:} 
$q_{1}=30,96$ [$\mu C$]$,$ $q_{2}=13,67$ [$\mu C$]$,$ $q_{3}=44,64$ [$\mu C$]%
\begin{equation*}
\FRAME{itbpFU}{1.7832in}{1.5869in}{0in}{\Qcb{Figura 2}}{}{Figure}{\special%
{language "Scientific Word";type "GRAPHIC";maintain-aspect-ratio
TRUE;display "USEDEF";valid_file "T";width 1.7832in;height 1.5869in;depth
0in;original-width 6.3771in;original-height 5.6723in;cropleft "0";croptop
"1";cropright "1";cropbottom "0";tempfilename
'R23O0G02.wmf';tempfile-properties "XPR";}}
\end{equation*}
\end{enumerate}

\item Para el circuito de la Fig. 3, $C_{1}=C_{5}=5,85$ [$\mu F$] y $%
C_{2}=C_{3}=C_{4}=6,70$ [$\mu F$]$.$ Si $V_{ab}=20,0$ [V] determine la
capacitancia equivalente y la carga que se acumula en cada condensador. 
\textbf{R: }$C_{eq}=2,26$ [$\mu F$], $q_{1}=q_{5}=45,2$ [$\mu C$], $%
q_{2}=30,55$ [$\mu C$]$,$ $q_{3}=q_{4}=15,28$ [$\mu C$] 
\begin{equation*}
\FRAME{itbpFU}{2.2079in}{1.4451in}{0in}{\Qcb{Figura 3}}{}{Figure}{\special%
{language "Scientific Word";type "GRAPHIC";maintain-aspect-ratio
TRUE;display "USEDEF";valid_file "T";width 2.2079in;height 1.4451in;depth
0in;original-width 7.3794in;original-height 4.817in;cropleft "0";croptop
"1";cropright "1";cropbottom "0";tempfilename
'R23O0G03.wmf';tempfile-properties "XPR";}}
\end{equation*}

\item Considere el condensador de placas paralelas de la Fig. 4, con $A=7,89$
$[cm^{2}]$, $d=4,62$ $[mm]$. Si la mitad superior est\'{a} llena con
polietileno ($\kappa _{2}=2,25$) y la mitad inferior con vidrio ($\kappa
_{1}=5,1$)$.$ Determine la capacitancia del dispositivo y la energ\'{\i}a
que puede almacenar con un voltaje $V=12,0$ $[V].$ \textbf{R: }$C_{eq}=4,71$
[pF] , $U=3,39\times 10^{-10}$ [J] 
\begin{equation*}
\FRAME{itbpFU}{1.7798in}{1.4313in}{0in}{\Qcb{Figura 4}}{}{Figure}{\special%
{language "Scientific Word";type "GRAPHIC";maintain-aspect-ratio
TRUE;display "USEDEF";valid_file "T";width 1.7798in;height 1.4313in;depth
0in;original-width 5.1249in;original-height 4.1182in;cropleft "0";croptop
"1";cropright "1";cropbottom "0";tempfilename
'R23O0H04.wmf';tempfile-properties "XPR";}}
\end{equation*}
\end{enumerate}

\end{document}
